\section{Related Work}\label{sec:related}

\scl{} is related to four lines of research: the trackers it works with, the
confidence of the detectors that feed them, trackers that adapt their own
thresholds, and methods that adapt inference while running. They are
reviewed below in this order, and Table~\ref{tab:related} compares the
closest methods with \scl{}.

Most practical trackers follow the tracking-by-detection design. SORT
predicts each box with a Kalman filter and matches the predictions to new
detections by overlap with the Hungarian method \cite{bewley2016sort}.
ByteTrack added a second matching step that uses low-score detections to
continue existing tracks \cite{zhang2022bytetrack}, and many later trackers
kept this two-stage design. Others, such as OC-SORT \cite{cao2023ocsort}, use
only one score threshold. All these trackers read detector scores through
fixed thresholds that were tuned for one detector, usually YOLOX
\cite{ge2021yolox}. Seven published trackers of both kinds are used in this
paper as hosts and are not changed (Sect.~\ref{sec:setup}).

These fixed thresholds are a weak point, because scores do not mean the same
thing for every detector. Detectors of different families, such as YOLOX
\cite{ge2021yolox}, YOLOv8 \cite{jocher2023yolov8}, YOLOv9
\cite{wang2024yolov9}, DETR \cite{carion2020detr}, and RT-DETR
\cite{zhao2024rtdetr}, give different score distributions. Their confidence
is not calibrated across models and conditions, and calibration methods for
object detection learn a mapping from labelled data
\cite{kuppers2020detcalib}. A tracker tuned for one detector therefore cannot
assume that another detector's scores mean the same thing, and labels to
recalibrate them are rarely available after deployment. Image processing
offers a tool that needs no labels: Otsu's method picks a threshold from a
histogram so that the two resulting groups are as separate as possible
\cite{otsu1979}. \scl{} applies it to the logits (log-odds) of the detector
scores of recent frames.

Some trackers already adapt their own thresholds. Adaptive ByteTrack sets the
score threshold of ByteTrack from the detections of each frame
\cite{ma2023adaptivebytetrack}. AdapTrack adapts the thresholds of its
matching step \cite{shim2024adaptrack}. ACT sets the high and low thresholds
of each frame from the first quartile of the scores and the number of
candidates \cite{act2026adaptive}. Other recent trackers use the score or the
spread of each detection inside the tracker. ConfTrack weights the Kalman
update, the Kalman noise, and the matching cost by the score of each box
\cite{jung2024conftrack}. UTrack measures the spread of the boxes that
non-maximum suppression (NMS) merges and uses it as the Kalman noise, so it needs the
boxes before suppression \cite{solano2025utrack}. PlugTrack blends the Kalman
filter with learned motion predictors through a small network trained with
labels \cite{kim2026plugtrack}. These methods are built into the tracker, use
only the current frame or each box alone, or need extra trained parts or the
inside of the detector. \scl{} differs in four ways. It uses past
frames. It decides for the whole stream whether to act at all, and otherwise
keeps the tracker's own settings. It works through a declared host contract
on several published trackers without changing them. And it is tested as one
frozen configuration under score changes and on external systems.

A wider group of methods adapts inference itself. Dynamic neural networks
adapt their computation to each input \cite{han2022dynamic}, and Chameleon
re-profiles the settings of video analytics pipelines over time
\cite{jiang2018chameleon}. These methods learn or measure a policy for each
model. Our earlier work, AC-MOT (a manuscript submitted for publication), is a controller that
needs no training and sets the detector operating point (score threshold,
NMS threshold, and input resolution) while running, from image and track
cues. Such a controller is tied to the detector it was calibrated with,
because a change in the detector's score scale does not change the image.
\scl{} reads the scores instead and acts on the tracker's input and
thresholds.

In summary, current methods either adapt thresholds or use the score of each
box inside one tracker, calibrate scores with labelled data, or learn a policy
for one model. What is missing is a layer that needs no training, adjusts itself
from the detector's own scores, works with many unchanged trackers, and keeps
a well-matched tracker at its own accuracy. \scl{} is designed to fill this
gap. Table~\ref{tab:related} summarizes the differences.

\begin{table*}[t]
\caption{Comparison of \scl{} with recent adaptive methods for tracking (BT: ByteTrack; ReID: re-identification network; NMS: non-maximum suppression).}
\label{tab:related}
\centering\footnotesize
\setlength{\tabcolsep}{4pt}
\begin{tabular}{@{}>{\raggedright\arraybackslash}p{24mm}c>{\raggedright\arraybackslash}p{30mm}>{\raggedright\arraybackslash}p{26mm}>{\raggedright\arraybackslash}p{19mm}>{\raggedright\arraybackslash}p{20mm}>{\raggedright\arraybackslash}p{28mm}@{}}
\toprule
Method & Year & Adapted & Statistics & Extra trained part & Detector input & Where it works\\
\midrule
Adapt.\ BT \cite{ma2023adaptivebytetrack} & 2023 & high/low score split & current frame & none & boxes and scores & inside BT\\
ConfTrack \cite{jung2024conftrack} & 2024 & Kalman update, noise, and cost & score of each box & ReID & boxes and scores & new tracker\\
AdapTrack \cite{shim2024adaptrack} & 2024 & matching thresholds & current frame & ReID & boxes and scores & new tracker\\
UTrack \cite{solano2025utrack} & 2025 & Kalman noise & spread of boxes merged by NMS & none & boxes before NMS & added to existing trackers\\
ACT \cite{act2026adaptive} & 2026 & high and low thresholds & current frame & none & boxes and scores & inside the tracker\\
PlugTrack \cite{kim2026plugtrack} & 2026 & mix of Kalman and learned motion & track boxes & network trained with labels & boxes & added to two learned predictors\\
K{\"u}ppers \cite{kuppers2020detcalib} & 2020 & detector scores & labelled data & calibration map & boxes and scores & detector output\\
Chameleon \cite{jiang2018chameleon} & 2018 & pipeline configuration & profiling & none & whole pipeline & analytics pipeline\\
AC-MOT (ours, submitted) & 2026 & detector operating point & image and track cues & none & sets its settings & before the detector\\
\scl{} (this paper) & -- & candidates, scores, thresholds & past frames, no labels & none & boxes and scores & between detector and seven unchanged trackers\\
\bottomrule
\end{tabular}
\end{table*}

